\section{Coefficients that detect analytic rank}
\label{cs:coefficient-systems}

The interpolation theorems require nonvanishing with a uniform valuation
bound at the nonzero vertices of a cube. We encode both requirements in
one binary coefficient: it is one when a modular coefficient has the
least normalized valuation. Waldspurger's formula supplies the even-sign
coefficients; reduction of genus Heegner sums supplies the odd-sign
coefficients at finite precision. We prove their modular and arithmetic
properties separately before choosing this normalization.

\subsection{Filters and prime weights}

Fix $E/\Q$ and a finite set $S$ containing the primes of $2N$ and the
ramification set of $F=\Q(E[2])$.  Additional fixed primes may be included.
A \emph{filter} consists of a sign $\sigma\in\{1,-1\}$ and unit
squareclasses at the primes in $S$, imposed on positive squarefree integers
$D$ prime to $S$, so that
\[
                     h=\sigma D
\]
is an odd fundamental discriminant and the functional sign of $E^{(h)}$
is constant.  Multiplication by squares of units preserves the local
conditions.  We impose these conditions on all Fourier indices, not only
on squarefree ones.  Projections onto the required unit squareclasses are
finite combinations of quadratic coefficient twists and coprimality
projections; they enlarge the level only at $S$.

For $p\notin S$, the action on $V=E[2]$ has the following three types:
\[
\begin{array}{c|c|c}
\text{type}&\Fr_p|_V&\text{weight}\\ \hline
\text{split}&1&1\\
\text{simple}&\text{an element of order }2&1/2\\
\text{root}&\text{an element of order }3&0.
\end{array}
\]
Thus $w(D)=\frac12\sum_{p\mid D}\dim_{\mathbf F_2}V^{\Fr_p}=w(h)$.
The fractional part of $w(D)$ is constant on a filter: it is determined
by the quadratic discriminant character of the residual representation.
For an arbitrary squarefree set of primes we use the same additive
weight.  Signs are encoded by $p^*=(-1/p)p$, so that $h=\prod_{p\mid D}p^*$.

We use the ordinary theta multiplier cubed in weight $3/2$. A finite
idele unit character and its Dirichlet character are related by the
reciprocal convention: if a finite Schwartz function obeys
$f(ux)=\alpha(u)f(x)$, its scalar classical character on the lower-right
entry is associated to $\alpha^{-1}$. All valuations satisfy $v(2)=1$.

\subsection{The even coefficient system}

\begin{proposition}[Even coefficients]
\label{prop:even-coefficients}\label{co:even}
Fix a nonempty filter with even functional sign.  There are a holomorphic
weight-$3/2$ cusp form
\[
                         g=\sum_{n\geq1}a(n)q^n
\]
of level supported on $S$, after enlarging exponents at those primes,
and finite-order characters $\psi,\eta$ of $2$-power order, with
$\psi/\eta$ quadratic, having the following properties.
\begin{enumerate}[label=\textup{(\roman*)}]
\item The coefficients vanish off the filter and, for $p\notin S$,
\begin{equation}
 a(p^2n)=\bigl(\eta(p)a_p(E)-\psi(p)(-n/p)\bigr)a(n)
                   -p\eta(p)^2a(n/p^2).
\label{cs:hecke-recurrence}
\end{equation}
As usual a coefficient with a nonintegral index is zero.
\item After normalizing at one nonzero squarefree index, every squarefree
coefficient in the filter is a rational number times a root of unity of
$2$-power order.  For a fixed nonzero squarefree index $D_0$,
\begin{equation}
 2v\left(\frac{a(D)}{a(D_0)}\right)
 =v\left(\frac{\mathcal L(\sigma D)}{\mathcal L(\sigma D_0)}\right)
 \quad\text{whenever }a(D)\ne0.
\label{cs:even-ratio}
\end{equation}
Equivalently, there is a constant $c_g$, independent of $D$, such that
\begin{equation}
       2v(a(D))=v\bigl({\cal L}(h)\bigr)+c_g.
\label{cs:even-value}
\end{equation}
In particular $a(D)\ne0$ if and only if $L(E^{(h)},1)\ne0$.
\item After the fixed auxiliary exclusions described below, there is a
constant $C_g$ such that, at every nonzero squarefree coefficient,
\begin{equation}
           v(a(D))\geq w(D)+\tfrac12s(h)-C_g.
\label{cs:even-lower-bound}
\end{equation}
The constants depend on the fixed curve, filter and auxiliary choices,
and not on $D$ or its number of prime factors.
\end{enumerate}
\end{proposition}

\begin{proof}
We first choose the packet member and its local characters. We then
establish the rationality up to phase needed for valuation comparisons,
apply the coefficient-square formula, and deduce the uniform lower bound
from the interpolation estimates.

Use the Shimura packet attached to the integral-weight newform of
$E^{(-\sigma)}$.  In the classical convention its positive Fourier index
$D$ tests the twist by $\chi_{-D}$, hence the curve $E^{(\sigma D)}$.
The packet and Fourier coefficient results of
\citet{Waldspurger1980,Waldspurger1981} apply to the cuspidal packet
orthogonal to unary theta series.  Their local uniqueness and global
multiplicity assertions are summarized in
\citep[Sections~3.1 and~5.1]{Waldspurger1983}.
A nonzero central value in the filter exists by
Proposition~\ref{prop:rough-nonvanishing}.  It selects a globally occurring
packet member with nonzero Fourier functional at that index.  At good
odd primes choose spherical vectors; at infinity choose the holomorphic
vector of weight $3/2$.

We explain the choices at a prime of $S$.  Let $\lambda_D$ be a nonzero
local Fourier functional and start with a vector on which it is nonzero.
Average upper-unipotent translates against the inverse additive character
of frequency $D$.  If the averaging compact subgroup is sufficiently
large, its translate by a unit scale $u$ can contribute only when $u^2$
is sufficiently close to $1$.  The surviving scales are consequently
small neighborhoods of $1$ and $-1$.  Local constancy then permits
projection onto a unit character which matches the forced action of
$-1$, without killing the Fourier value: the contributions of the two
neighborhoods add.  Such a character may be chosen to have $2$-power
order.  Indeed only its parity is prescribed, and the $2$-primary
quotient of the unit group realizes either possible parity.  Upper
averaging at integral frequency gives period one, and lower-unipotent
invariance holds at sufficiently large depth.  These vectors therefore
give a classical form at a level supported on $S$.  The parity conditions
combine correctly by automorphy.  The subsequent squareclass projections
preserve its character.  Write that character as $\psi$.

In the theta-multiplier convention the good eigenvalues of the packet
are $\psi(p)$ times those of $E^{(-\sigma)}$.  Set
$\eta=\psi\chi_{-\sigma}$.  Then $\psi/\eta$ is quadratic and the
weight-$3/2$ Hecke formula gives Equation~\eqref{cs:hecke-recurrence}.

We need more than algebraicity of the coefficients. After normalization
at one squarefree index, their nonzero squarefree valuations must be
integers, so that a strict valuation inequality gains a full factor of
$2$, even when the character field is ramified. We prove this by showing
that every such coefficient is rational up to a phase of $2$-power order.
At the fixed level,
the weight-$3/2$ cusp space with character $\psi$ has a basis over the
field of character values, and coefficient conjugation changes $\psi$
to its conjugate
\citep[Section~2, Lemma~4(3)--(4), pp.~658--659]{Shimura1981CriticalValues}.
The good Hecke equations and the vanishing of coefficients off the
filter are linear equations over $\Q(\psi,\eta)$.  In a finite-dimensional
space, finitely many of these equations cut out their full common
kernel.  The constructed complex vector has a nonzero coefficient at
a squarefree index $D_0$ of the filter, so this coefficient functional
is nonzero on some algebraic vector in that kernel.  Choose such a
vector.  It remains orthogonal to unary theta series: at good primes
their eigenvalues have absolute value $p+1$, whereas the prescribed
eigenvalues have absolute value at most $2\sqrt p$ by Hasse's bound
\citep[Section~2, Lemma~2.1, and Corollary~5.2]{Purkait2012}.

We use \citet[Corollary~2, pp.~483--484]{Waldspurger1981} to compare
Fourier coefficients.  Its hypothesis (H1) holds for every nonzero
classical eigenspace used here, by Proposition~2 of that paper.
The final paragraph of the corollary's proof removes (H2) when the
indices lie in the same local squareclass at every level prime.
Our filters retain this condition at~$2$ as well as at the odd primes.

Let $u$ be an automorphism of $\overline\Q$ acting by an odd power $e$
on the $2$-power characters, and put $\nu=\eta^{(e-1)/2}$.  The
coefficient conjugate $g^u$ and coefficient twist $g\otimes\nu$ have
exactly the same character and good Hecke system, since
\[
 \psi\nu^2=\psi\eta^{e-1}=\psi^e,
                 \qquad \eta\nu^2=\eta^e.
\]
They may be compared at a common level supported on $S$, in the same
good Hecke eigenspace $W$, and both have a nonzero $D_0$-coefficient.
Let $\ell_n:W\to\C$ denote the $n$th Fourier-coefficient functional.
Waldspurger's corollary applies to every $f\in W$ and gives
\[
                    \ell_D(f)^2=C_D\ell_{D_0}(f)^2,
\]
where $C_D$ is independent of $f$.  If $C_D=0$, then $\ell_D=0$.
Otherwise choose $c_D\in\C$ with $c_D^2=C_D$ and factor the identity
as
\[
       (\ell_D-c_D\ell_{D_0})(\ell_D+c_D\ell_{D_0})=0.
\]
The polynomial ring on $W$ is an integral domain, so one factor
vanishes identically.  Thus the two coefficient functionals are
proportional on the whole space, including oldvectors and linear
combinations.  In particular the two forms have the same ratio of
their $D$- and $D_0$-coefficients.

It follows, for $\xi=a(D)/a(D_0)$, that
\[
                         u(\xi)=\nu(D/D_0)\xi.
\]
The local squareclass conditions imply
$\eta(D/D_0)=z^2$ for a value $z$ in the group of values of $\eta$.
Consequently $u(\xi/z)=\xi/z$ for every $u$, so $\xi/z\in\Q$.
This proves the phase assertion.

With our characters and period normalization, its identity of
\emph{coefficient squares} is
\begin{equation}
 \left(\frac{a(D)}{a(D_0)}\right)^2
 =\psi(D/D_0)\,
   \frac{\mathcal L(\sigma D)}{\mathcal L(\sigma D_0)}.
 \label{cs:even-square-ratio}
\end{equation}
Here $\psi(D/D_0)=\psi(D)\psi(D_0)^{-1}$ has valuation zero;
the factor $\sqrt{D/D_0}$ in the classical central-value comparison
is incorporated in $\mathcal L$. The local unit
squareclasses at $S$ agree; at a good prime the spherical squarefree
factor is one at both valuations zero and one
\citep[Section~I.3, pp.~377--378]{Waldspurger1981}.
There is consequently no extra product over the prime factors of $D$.
Taking valuations gives Equation~\eqref{cs:even-ratio}, hence
Equation~\eqref{cs:even-value}. The same identity, compared with the
nonzero coefficient at $D_0$, gives the asserted equivalence of
vanishing. The corresponding rational-square statement for central values
is Vign\'eras's rationality theorem, recalled in
\citet[Section~6.3, Theorem~6]{Waldspurger1983}.

If $E[2]$ is an extension of trivial modules, combine
Equation~\eqref{cs:even-value} with
Theorem~\ref{thm:cyclotomic-bound}.
If $E[2]$ is irreducible, choose a fixed odd fundamental discriminant
$h_*$ of opposite infinite sign, of analytic rank one, whose local
characters agree with those of $h$ at $2N$.  Proposition~\ref{prop:rough-nonvanishing}
supplies such a choice.  Exclude its prime factors from the filter and
choose its support to avoid the finitely many exceptional genus fields.
The reverse genus comparison in Proposition~\ref{prop:heegner-comparison}
and Theorem~\ref{thm:heegner-genus-bound} give the same lower bound.
These are fixed exclusions and constants, proving
Equation~\eqref{cs:even-lower-bound} in both cases.
\end{proof}

\subsection{Odd coefficients from reduction of Heegner traces}

The odd construction combines reduction of Heegner points with the ternary
quadratic forms attached to definite quaternion orders. We use Deuring's
reduction correspondence in the oriented form of Cornut--Jetchev,
together with the Eichler-order and ternary-form description of
Jetchev--Kane \citep{Deuring1941,CornutJetchev2013,JetchevKane2011}.
The weights below retain the genus character and the local orientations,
so that the resulting Fourier coefficients recover the reduced genus sum
at finite precision.

Fix now a nonempty filter of odd functional sign.  Choose an odd fundamental
$h_*$ of opposite infinite sign, with $L(E^{(h_*)},1)\ne0$, whose local
characters agree with those of $h$ at $2N$.  Exclude its primes.  We may
choose $|h_*|>3$ and include in its support a prime outside the
exceptional imaginary quadratic discriminants.  Thus all genus fields
below have only the units $\{1,-1\}$ and avoid the CM field of $E$.
Write
\[
          k=hh_*<0,\qquad K_D=\Q(\sqrt{k}),\qquad
          P_D=P_1(h_*).
\]
The point $P_D$ is the character \emph{sum} of the conductor-one Heegner
point for $K_D$. We fix the modular parametrization $\varphi$ with
$\varphi(\infty)=0$, so the constants in its Hecke and Atkin--Lehner
relations are torsion.  At a simple zero of $E^{(h)}$, its free line is that of
$E^{(h)}$, up to the fixed-degree identifications in
Proposition~\ref{prop:heegner-comparison}.
Below, the product derivative means the derivative at $s=1$ of
$L(E^{(h)},s)L(E^{(h_*)},s)$.

\begin{proposition}[Odd coefficients under a genus-index bound]
\label{prop:odd-coefficients}\label{co:odd}
Assume that a constant $C_*$ satisfies
\begin{equation}
              j(P_D)\geq w(D)+\tfrac12s(h)-C_*
\label{cs:odd-bound-hypothesis}
\end{equation}
for every $h$ in the filter with a simple zero.  There are arbitrarily
large precisions $M$, good odd primes $\ell$ outside the fixed support,
and holomorphic weight-$3/2$ forms $g_M=\sum a_M(n)q^n$, of level
supported on $S\cup\{\ell\}$, with the following properties.
\begin{enumerate}[label=\textup{(\roman*)}]
\item The additional filter at $\ell$ is
\begin{equation}
  (h_*/\ell)=-\sigma,\qquad (h/\ell)=\sigma,
               \qquad \rho_T(\Fr_\ell)\equiv C\pmod{2^M},
\label{cs:detecting-prime}
\end{equation}
where $C$ is the action of a complex conjugation.  In particular
$\ell\equiv-1\pmod{2^M}$.
\item After a fixed harmless decrease in the working precision, the
forms satisfy Equation~\eqref{cs:hecke-recurrence} modulo that precision,
with characters $\eta_M,\psi_M$ of $2$-power order and quadratic ratio.
Their squarefree coefficients have the form of a $2$-power phase times
an integer modulo $2^M$.
\item If a coefficient is visible, meaning $v(a_M(D))<M$, then
$E^{(h)}$ has a simple zero and
\begin{equation}
                    j(P_D)\leq v(a_M(D))+C_1.
\label{cs:odd-detection}
\end{equation}
For every admissible $D$ the coefficient lower bound is
\begin{equation}
 v(a_M(D))\geq
       \min\{M,\,w(D)+\tfrac12s(h)-C_2\}
\label{cs:odd-lower-bound}
\end{equation}
when the product derivative is nonzero; if it is zero, the coefficient
is zero modulo $2^M$.  Here $C_1,C_2$ are independent of $\ell,M,D$.
\item A fixed squarefree witness $D_w$ in the original filter is
admissible at every chosen $\ell$ and has $v(a_M(D_w))$ bounded as
$M\to\infty$.
\end{enumerate}
For irreducible $E[2]$, the hypothesis is
Theorem~\ref{thm:heegner-genus-bound}.  For rational two-torsion, the positive filters used in the pointwise
argument have this bound by Proposition~\ref{prop:rational-odd-bound},
using only the even source. An arbitrary twist is first absorbed into
the fixed curve, so only these positive filters are needed there.
\end{proposition}

\begin{proof}
The proof has three stages. A reduction map detects one fixed nonzero
genus point at arbitrarily high precisions. Weighted theta series then
supply holomorphic forms with the required Hecke recurrence. Finally,
an exact comparison with genus sums transfers point divisibility and
the assumed lower bound to their squarefree coefficients.

\paragraph{The reduction map and a detecting sequence.}
Choose a nonzero integral row $b$ with $bC=\sigma b$.  Under
Equation~\eqref{cs:detecting-prime}, finite-field Kummer theory gives
\[
 \kappa_M:E(\mathbf F_{\ell^2})\longrightarrow
 H^1(\mathbf F_{\ell^2},E[2^M])
       =E[2^M]/(\Fr_\ell^2-1)E[2^M]=E[2^M].
\]
The last identification is evaluation at arithmetic Frobenius squared.
Put $\lambda_M=b\kappa_M$.  This map has Frobenius sign $\sigma$.

Choose $h_w=\sigma D_w$ in the original filter with a simple zero,
using Proposition~\ref{prop:rough-nonvanishing}.  A fixed multiple of
its genus point is a nonzero point on the rational curve $E^{(h_w)}$.
Its division cocycle has open translation image on the kernel of the
Tate action, even after fixed finite auxiliary fields have been
adjoined.  Here is the relevant argument.  The translation image is a
Galois-stable $\Ztwo$-submodule of the Tate module.  Its rational span is
nonzero: otherwise the point's Kummer class would factor through the
Tate image; that image has a central nonidentity scalar, which kills
its rational first cohomology, contradicting Kummer injectivity.
The rational Tate representation is irreducible by
\citet[Chapter~IV, Section~2.1, pp.~IV-9--IV-10]{Serre1968}, since
$\operatorname{End}_{\Q}(E)=\Z$.  This includes curves with geometric
complex multiplication, whose nonreal endomorphisms are not defined
over $\Q$.  The scalar Lie-algebra assertion of
\citet[Section~2, Theorem~3, p.~62]{Bogomolov1981}, followed by local
logarithm and exponential, supplies the required nonidentity scalar.
The nonzero stable span is therefore the whole Tate module.  The
translation image consequently contains a full lattice.  Passing to
a finite-index subgroup preserves this fact.

Take an element acting as $C$ on $T$ and as complex conjugation on the
fixed quadratic character fields, and modify it by such a translation.
On the twisted module for $h_w$ its action is $\sigma C$; at its square
the cocycle includes $1+\sigma C$.  Projection by $b$ is twice the
corresponding nonzero row.  We can therefore arrange nonzero detection
with a fixed finite valuation.  Chebotarev, applied at successive finite
precisions to this fixed compatible Galois element and its division
data, supplies arbitrarily large primes satisfying
Equation~\eqref{cs:detecting-prime}, also with $h=h_w$, and the asserted
bounded detection.  A congruence on the target curve's Frobenius is all
that is required.  The supersingular objects used next are points of
the modular curve, not an additional reduction hypothesis on $E$.

\paragraph{Supersingular values and local weights.}
Let $B$ be the quaternion algebra ramified at $\ell$ and infinity, and
let $R$ be an Eichler order of level $N$.  Supersingular pairs on $X_0(N)$
are parametrized by
\[
                  B^\times\backslash\widehat B^\times/
                  \widehat R^\times.
\]
They are defined over $\mathbf F_{\ell^2}$.  The reduction and Eichler
order description apply to arbitrary $N$ prime to $\ell$; see
\citep[Sections~1.1--1.2 and~2.5]{JetchevKane2011}.
For a representative $z$ write $R_z=B\cap z\widehat Rz^{-1}$ and let
$P_M(z)$ be the reduced modular-parametrization value followed by
$\lambda_M$.  The cusp differences have torsion image by
\citet[Section~1]{Drinfeld1973}.  Multiply the parametrization once by a fixed integer
$L$ killing the cuspidal constants in its Hecke and Atkin--Lehner
relations and the bounded-degree torsion occurring in genus sums.
This choice is independent of $D,\ell,M$.

The resulting $P_M$ has good Brandt eigenvalues $a_p(E)$, level
Atkin--Lehner eigenvalues $w_{p^e}$ for $p^e\Vert N$, and eigenvalue
$\sigma$ for the right uniformizer at $\ell$.  These identities hold
modulo $2^M$.  Choose integer representatives of its values when making
an ordinary characteristic-zero theta series.

On $B^0$, the trace-zero subspace, use the positive definite quadratic
form $Q(x)=-x^2=\Nrd(x)$. Our aim is to make the squarefree coefficient
$a_M(D)$ equal, modulo $2^M$ and up to a phase of $2$-power order, to a
fixed integer multiple of
\[
                 \lambda_M\bigl(\operatorname{red}_{\ell}(L P_D)\bigr).
\]
We will identify its terms with optimal embeddings of
$\mathcal O_{K_D}$. Their local orientations give several copies of the genus sum.
The weights must make those copies contribute equally, while their
unit-scaling law must give a scalar modular character. The weighted
trace calculation below will determine the exact multiplier.

Define a finite compactly supported, locally constant weight
$f=\prod f_p$ as follows; a value not satisfying a listed condition is zero.
\begin{enumerate}[label=\textup{(\alph*)}]
\item Away from $2Nh_*\ell$, require order integrality and put $f_p=1$.
\item At $2$, require $(1+x)/2$ to be integral and $x^2$ to be a split
fundamental unit squareclass.  At $p^e\Vert N$, require the split
fundamental unit type and its order embedding.  The cyclic Eichler
level line selects one of the two eigenroots $r$ of $x$; weight it by a
$2$-power unit character $\eta_p(r)$ with
\[
                   \eta_p(-1)=w_{p^e}\chi_{h_*}(p)^e.
\]
If $p=2$, use the selected summand of $(1+x)/2$, whose eigenvalues
$(1+r)/2$ and $(1-r)/2$ differ by the odd unit $r$. Scaling by an odd
$u$ sends this operator to
$u(1+x)/2+(1-u)/2$ and preserves that summand; its selected root for
$ux$ is $ur$. Thus the weight changes by $\eta_2(u)$, even when the
two eigenvalue parities interchange. If $2\nmid N$, impose only the
stated integrality restriction.
\item At $p\mid h_*$ require integral $x$ with determinant of valuation
one.  Its reduction is nonzero nilpotent.  On that orbit use the sign
which transforms by $(\det u/p)$ under conjugation by
$u\in\GL_2(\mathbf Z_p)$.  This sign exists because the stabilizer of
a nonzero nilpotent has the form $aI+bX$ and determinant $a^2$.
Scaling $x$ by a unit $u$ changes this sign by $(u/p)$.
\item At $\ell$ require integral unit norm and the unramified quadratic
type.  Choose a basis of the trace-zero line in $\mathbf F_{\ell^2}$
and weight the nonzero residue on that line by its Legendre symbol
relative to the basis.
\end{enumerate}
All these conditions may be imposed at finite depth.  They are stable
under unit scaling.  Right order-unit conjugation changes $f$ by the
inverse of $\chi_{h_*}\circ\Nrd$, exactly canceling the norm character
below.  At $\ell$, order units commute on the residue field, while a
uniformizer negates its trace-zero line.  Since $\ell\equiv-1\pmod4$,
the latter weight changes by $-1$.

Write $\alpha_B=\prod_p\alpha_p$ for the finite unit-scaling
character, so that $f(ux)=\alpha_B(u)f(x)$.  Define the classical
Dirichlet character $\eta_B$ by
\[
 \eta_B(d)=\prod_p\alpha_p(d)
 \qquad\text{for integers }d\text{ prime to the support}.
\]
The associated finite idele character restricts to $\alpha_B^{-1}$
on the ramified unit groups, in accordance with our reciprocal
convention.  Both characters are even.  Indeed their value on $-1$ is
\[
 \left(\prod_{p^e\Vert N}w_{p^e}\chi_{h_*}(p)^e\right)
                      \chi_{h_*}(-1)(-1)=1,
\]
by the even functional sign of $E^{(h_*)}$.  The characters have
$2$-power order.

\paragraph{Transformation and good Hecke relation.}
Choose a fixed positive integer $c_0$ divisible by the orders of all
unit groups which occur.  Such a choice is uniform: a unit group of a
rational definite quaternion order has bounded order.  Form
\begin{equation}
 H_M(n)=c_0\sum_z\frac{P_M(z)\chi_{h_*}(\Nrd z)}{|R_z^\times|}
             \sum_{\substack{x\in B^0\\-x^2=n}}f(z^{-1}xz),
 \qquad H_M(\tau)=\sum_{n\geq1}H_M(n)q^n.
\label{cs:weighted-theta}
\end{equation}
The norm character is trivial on rational $B^\times$, since its rational
reduced norms are positive.  The right transformation just checked makes
Equation~\eqref{cs:weighted-theta} independent of the representatives.

For scalar modularity, apply the Weil representation of $(B^0,Q)$
to $f$ and the Gaussian at infinity \citep{Weil1964}.  At a finite prime,
with the additive character of conductor $\Z_p$, the local formulas are
\citep[Section~4.3, pp.~1965--1966]{PalDeVeraPiquero2019}
\[
 \begin{aligned}
 \omega(n(b))f_p(x)&=\mathbf e_p(bQ(x))f_p(x),\\
 \omega(m(u))f_p(x)&=\gamma_p(u)|u|_p^{3/2}f_p(ux),
 &\omega(w)f_p&=\gamma_p\widehat f_p.
 \end{aligned}
\]
Here $n(b)=\left(\begin{smallmatrix}1&b\\0&1\end{smallmatrix}\right)$,
$m(u)=\operatorname{diag}(u,u^{-1})$, and $w$ is the Weyl element.
Integral $b$ fixes $f_p$, since $Q$ is integral on its support.
Choose a lattice of translation invariance for $f_p$; its dual contains
the support of $\widehat f_p$. For sufficiently large $a_p$, the form
$p^{a_p}Q$ is integral on that dual lattice. Thus
$n(p^{a_p}b)$ fixes $\widehat f_p$, and conjugation by $w$ proves
invariance under sufficiently deep lower unipotents. This argument
also works at $2$ and at arbitrary finite character depth.

Compare with the threefold ordinary-theta reference line, using the
same additive character and lifts.  In a presentation $B=(-a,-b)$,
the quadratic form is $aX^2+bY^2+abZ^2$.  The ratio of its diagonal
Weil factor to the reference factor is
\[
                   (a,u)_p(b,u)_p(ab,u)_p=((ab)^2,u)_p=1.
\]
This calculation holds also over $\Q_2$.  The remaining local
diagonal scalar is therefore $\alpha_p(u)$.  Upper unipotents,
diagonal units, and sufficiently deep lower unipotents generate the
required local $\Gamma_0$ subgroup: explicitly, for $a$ a unit,
\[
 \begin{pmatrix}a&b\\c&d\end{pmatrix}
 =n^-(c/a)m(a)n(b/a).
\]
Choose the depths also to serve the ordinary-theta reference line
and the character conductors.  The central ambiguity of the
metaplectic lifts cancels in the comparison, since both representations
have the same odd central action.  At the chosen depth, $a$ is congruent
to $d^{-1}$ modulo the character conductor, so the local right
eigenvalue is $\alpha_p(d)^{-1}$.  Classical
restriction uses the inverse finite component and consequently gives
the character $\eta_B(d)$.  Conjugation by $z$ preserves $Q$ and
commutes with the Weil action, so all summands have this same level
and scalar character.
At good odd primes the lattice is spherical. The Gaussian gives
weight $3/2$, and positive definiteness gives holomorphy at every cusp.
Thus $H_M$ is a classical form of level supported on $2Nh_*\ell$.

At a prime $p\nmid 2Nh_*\ell$, the neighbor relation has Hecke eigenvalue
\[
                         \eta_B(p)\chi_{h_*}(p)a_p(E).
\]
The factor
$\chi_{h_*}(p)$ comes from the reduced norm of a Brandt neighbor and
$\eta_B(p)$ from scaling at the other primes.

The multiplicities can be read directly in a matrix order.  Put $y=px$.
If $y$ is primitive
and $Q(y)$ is divisible by $p^2$, its nonzero nilpotent reduction has a
unique invariant line; in a basis adapted to that line the other
off-diagonal entry is divisible by $p^2$, so exactly one neighbor makes
$y/p$ integral.  If $y/p$ is integral, its reduction has
$1+(-Q(y/p)/p)$ invariant lines, except that a zero reduction has a
further $p$ lines.  These are precisely the middle and last terms of
the weight-$3/2$ Hecke recurrence.  The factors $|R_z^\times|^{-1}$
are the orbit-stabilizer factors making the Brandt transfer valid.
Since the eigenrelation for $P_M$ holds modulo $2^M$, so does this
relation for $H_M$.

Put $d=|h_*|$. The weights force $Q(x)$ to be divisible by $d$.
Before the final filter projections, the series with coefficients
$a_M(n)=H_M(dn)$ is the same weighted theta sum for $Q/d$.
Its ternary determinant changes by $d^{-3}$, whose squareclass is $d$;
hence the scalar character acquires $(d/\cdot)$.
Substitution of $dn$ in the good-prime recurrence leaves
$\eta_B\chi_{h_*}$ unchanged. Now apply the filter projections.
Extraction of these multiples changes the character to
$\psi_M=\eta_B( |h_*|/\cdot)$, while
$\eta_M=\eta_B\chi_{h_*}$ is the good eigenvalue character.
Their quotient is quadratic.  The neighbor relation therefore becomes
Equation~\eqref{cs:hecke-recurrence}.  Every operation has level supported
on $S\cup\{\ell\}$.

\paragraph{The weighted trace identity.}
The form and its recurrence are now constructed. We next identify a
squarefree coefficient with the reduction of the original genus sum,
keeping the orientation multiplicities and unit stabilizers explicit.

For a squarefree $D$ in the filter, put $n=|h_*|D=-k$.
An $x$ counted in Equation~\eqref{cs:weighted-theta} is exactly an
optimal embedding of the full order $\mathcal O_{K_D}$ into $R_z$,
by sending $\sqrt{k}$ to $x$.  At $2$ the generator is
$(1+\sqrt{k})/2$, which is why the integrality condition above is
necessary.  At a split level prime there are two choices of cyclic
level orientation, and at $\ell$ there are two choices of unramified
quadratic orientation.  There is one local orbit at the remaining
primes.  The split-level assertion follows by decomposing a stable
lattice into its two idempotent summands: a cyclic quotient of order
$p^e$ forces all of that quotient into one summand.  The division-prime
assertion follows from the two reduced-norm parity cosets.  At primes outside $2Nh_*\ell$,
stable lattices are the fractional ideals of the quadratic
order, giving the single local orbit.

Fix one orientation tuple and a rational embedding $x_0$.  The
corresponding global optimal-embedding orbits are parametrized by
$\Pic(\mathcal O_{K_D})$.  Explicitly, once one adelic representative
$z_0$ is fixed, the representatives are $z_{\mathfrak a}=a z_0$ with
$a\in\widehat K_D^\times$; changing $a$ by $K_D^\times$ or
$\widehat{\mathcal O}_{K_D}^{\times}$ gives exactly the same class.
Because $a$ commutes with $x_0$, the Schwartz factor is constant along
this class-group orbit.  Its remaining relative weight is
\[
             \chi_{h_*}(\Nrd a)
              =\chi_{h_*}(\operatorname{Norm}_{K_D/\Q}a),
\]
the unramified genus character.  This computation establishes the
weighted identity, rather than merely an unweighted correspondence.

For completeness the stabilizer factor is also exact.  Each orbit of
embeddings in a given $R_z$ has
$|R_z^\times|/|\mathcal O_{K_D}^\times|=|R_z^\times|/2$ elements.
Thus its contribution after division by $|R_z^\times|$ is one half of
the corresponding class-group term. These are oriented embeddings,
not conjugate pairs: $x$ and $-x$ have opposite orientations at $\ell$
and cannot be conjugate by an order unit, whose residue conjugation
on $\mathbf F_{\ell^2}$ is trivial.  CM reduction at the inert prime
identifies these terms with the reductions of the Heegner points and
commutes with ideal action; the cyclic subgroup carries the selected
level orientation.  Fixing the residue-field identification at $\ell$
retains its orientation in the reduction-to-embedding correspondence
of \citet{Deuring1941}; see
\citep[arXiv version~1, Theorem~3.1, Corollary~3.2, and Section~3.6,
Remark~1]{CornutJetchev2013} and
\citep[Sections~2.4--2.6 and~4.1]{JetchevKane2011}.
The former remark explains that forgetting this identification
identifies conjugate embeddings.  Here both orientations remain in
the theta sum.

Choose the initial orientation to agree with the Heegner point defining
$P_D$.  The constant weight on its orbit is a phase
$\zeta_M(D)\in\mu_{2^\infty}$, since all its local weights are units.
Changing orientation at $p^e\Vert N$ multiplies the Schwartz factor by
$w_{p^e}\chi_{h_*}(p)^e$, and multiplies the norm-character-weighted
parametrization value by the same sign.  At $\ell$ both signs are $-1$,
since $\sigma\chi_{h_*}(\ell)=-1$.  Hence every one of the
$2^{\omega(N)+1}$ orientation tuples contributes the same term.
With the fixed multiplier $L$ already incorporated into $P_M$, the
comparison is consequently
\begin{equation}
 a_M(D)=c_0\,2^{\omega(N)}\zeta_M(D)\,
           \lambda_M\bigl(\operatorname{red}_{\ell}(L P_D)\bigr)
                           \pmod{2^M}.
\label{cs:weighted-trace-identity}
\end{equation}
The initial basis at $\ell$ and initial orientation affect only
$\zeta_M(D)$.  Formula~\eqref{cs:weighted-trace-identity} specifies the
normalization using a character sum, with no class-number denominator.
All rational multipliers are fixed independently of $D$ and $\ell$.

If the product derivative is zero, Gross--Zagier makes $P_D$ torsion.
The fixed multiplier kills it: genus sums and their twist comparisons
are defined over biquadratic fields.  The torsion of a fixed elliptic
curve in degree at most four is uniformly bounded, as one sees by
reducing its prime-to-residue-characteristic torsion at two fixed good
primes of distinct residue characteristic; the residue degrees are at
most four.  Thus the coefficient is zero at
the stated precision.  If the derivative is nonzero, write the point
against the free twist line.  Division by $2^{j(P_D)}$ on that line
implies the same divisibility after reduction and $\lambda_M$, up to
the fixed-degree comparison and the fixed multiplier.  Conversely,
a visible image bounds $j(P_D)$ above by its visible valuation plus a
fixed constant.  Equation~\eqref{cs:weighted-trace-identity} proves
Equations~\eqref{cs:odd-detection} and~\eqref{cs:odd-lower-bound}, using
Equation~\eqref{cs:odd-bound-hypothesis}.  Gross--Zagier--Kolyvagin
then gives the simple zero and finite $\Sha$ when the coefficient is
visible.  The detecting construction for $D_w$ supplies the final
assertion.  Fixed denominators or multipliers only reduce $M$ by a
fixed amount; reindexing the precisions gives the formulation stated.
\end{proof}

\subsection{The least normalized depth}

We now minimize the coefficient valuation after subtracting the prime
weight $w(D)$. This normalization is fixed before any prime array is
chosen. Attainment of the minimum will detect analytic rank and bound
the finite Selmer contribution; for the odd source this is interpreted
through the bounded-weight sequences and finite realizations below.

For the even source, first rescale $g$ so that $a(D_0)=1$ at one
nonzero squarefree index. Its nonzero squarefree valuations are then
integers by Proposition~\ref{prop:even-coefficients}. Set
\begin{equation}
              \delta=\min_D\{v(a(D))-w(D)\},\qquad
 A(D)=\begin{cases}
 1,&v(a(D))=\delta+w(D),\\
 0,&\text{otherwise},
 \end{cases}
\label{cs:even-symbol}
\end{equation}
with $A(D)=0$ off the filter.  The minimum exists: the lower bound is
uniform, visible valuations are integers after normalization, the
fractional part of the weight is fixed, and a nonzero witness exists.

For the odd source fix a nonprincipal ultrafilter $\mathcal U$ on its
sequence of precisions. Consider all sequences $D_M$ of admissible indices for
which $w(D_M)$ is bounded and $v(a_M(D_M))$ is visible on an
ultrafilter-large set.  For each such sequence, its normalized depth
has a lower bound in a discrete set.  Let $\delta$ be the least depth
attainable on an ultrafilter-large set.  This least depth exists and is
attained: the uniform lower bound bounds it below, discreteness gives
a least value, and the fixed witness bounds the possible minimum above.
For a bounded-weight admissible sequence $\mathbf D=(D_M)$, define
\[
 A(\mathbf D)=1
 \quad\Longleftrightarrow\quad
 \{M:v(a_M(D_M))=\delta+w(D_M)\}\in\mathcal U,
\]
and set $A(\mathbf D)=0$ otherwise. We suppress the sequence notation
when discussing finite configurations; each application uses their
simultaneous finite realizations, as proved below.

For the odd source, a least weight supporting a unit symbol is minimized
among the same bounded-weight sequences that attain depth $\delta$.
After restricting to an ultrafilter-large set, the weight of such a
sequence may be held constant, since it lies in a bounded subset of
$\tfrac12\Z_{\ge0}$. The set of these constant weights is nonempty and
has a least element $m$. If unit symbols of weight below $m$ occurred
on an ultrafilter-large set of stages, choosing one such index at each
stage would contradict the definition of $m$.

\begin{lemma}[Uniformity at a fixed weight bound]
\label{cs:fixed-weight-uniformity}
For every fixed $W$, on an ultrafilter-large set of sufficiently high
precisions,
\[
                   v(a_M(D))\geq\delta+w(D)
\]
for every admissible index with $w(D)\leq W$, with valuations truncated
at the working precision.  Any finite list of stabilized symbol
assertions of bounded weight is realized simultaneously at such
precisions.  The definition allows arbitrarily many weight-zero root
factors.
\end{lemma}

\begin{proof}
If the first assertion failed on an ultrafilter-large set, choose a
violating $D_M$ at each of those levels.  Their weights are bounded by
$W$ and their normalized depths lie in a fixed bounded discrete range.
After selecting one depth on an ultrafilter-large set, they contradict
the definition of $\delta$.  For a finite collection of stabilized
bits, intersect their ultrafilter-large realization sets.  A bound on
weight, rather than on the number or size of primes, is all that this
argument uses.  At weights beyond the chosen bound, the truncated
lower estimate of Proposition~\ref{prop:odd-coefficients} still
supplies any prescribed bounded-depth divisibility test once the weight
is sufficiently large.  No claim of one finite level valid for all
unbounded weights is made.
\end{proof}

\begin{proposition}[Detection by a unit symbol]
\label{prop:unit-symbol-detection}
For either coefficient source, $A(D)=1$ implies the prescribed analytic
order: zero in the even case and one in the odd case.  At every such
index,
\begin{align}
 v({\cal L}(h))&\leq2w(D)+C &&\text{in the even case},
 \label{cs:even-sharp}\\
 j(P_D)&\leq w(D)+C &&\text{in the odd case},
 \label{cs:odd-sharp}\\
 s(h)&\leq C,\qquad \dim_{\mathbf F_2}\Sel_2(E^{(h)}/\Q)\leq C.
 \label{cs:unit-selmer-bound}
\end{align}
The constants are fixed after choosing the original curve, filter,
coefficient normalization and genus partner.  For the odd source they
are independent of the detecting prime and the working precision, and
of the number and size of active primes.  Statements about its symbols
are used through finite realizations as in
Lemma~\ref{cs:fixed-weight-uniformity}.
\end{proposition}

\begin{proof}
In the even case substitute $v(a(D))=\delta+w(D)$ into
Equations~\eqref{cs:even-value} and~\eqref{cs:even-lower-bound}.
In the odd case use a precision greater than the finitely many tested
depths and apply Equations~\eqref{cs:odd-detection} and
\eqref{cs:odd-lower-bound}.  The constant $\delta$ is fixed and the
comparison constants are uniform, proving the first inequalities and
the bound on $s(h)$.  At these indices the analytic order is established,
so Gross--Zagier--Kolyvagin gives the same algebraic rank and finite
$\Sha$ \citep{GrossZagier1986,Kolyvagin1990}.  The exact sequence for $2$-Selmer then bounds its dimension
by that rank, the fixed bound on rational $2$-torsion, and the
$2$-primary length of $\Sha$.  This proves the final assertion without
making any finiteness assumption at an index where the analytic order
has not yet been detected.
\end{proof}
